% Part: proof-theory
% Chapter: sequent-calculus
% Section: introduction

\documentclass[../../../include/open-logic-section]{subfiles}

\begin{document}

\olfileid{pt}{seq}{int}

\olsection{引言}

相继式演算是由 Gerhard Gentzen（根岑）在二十世纪三十年代引入的推导系统。与其说它们的设计初衷是一种实用的推导方法，不如说 Gentzen 是将它们作为一种工具，以便能够证明有关推导的可能结构及其性质的某些结果。这类问题正是证明论所关注的。

正如其名，相继式演算处理的对象是\emph{相继式}，而非公式。一个相继式是由公式序列或公式多重集组成的对。Gentzen 最初的系统（\Log{LK} 和 \Log{LJ}）使用序列；现在的证明论学者则更偏好使用多重集。多重集类似于集合，不同之处在于一个多重集可以多次包含同一个元素。但与集合一样，元素的顺序无关紧要。因此，表达式 $\{!A, !A, !B\}$ 和 $\{!A, !B, !A\}$ 描述的是同一个多重集，但该多重集不同于 $\{!A, !B\}$ 和 $\{!A, !A, !B, !B\}$，因为后两者包含的 $!A$ 和 $!B$ 的数量与前者不同。

传统上，证明论中的公式多重集用大写希腊字母表示。所以当我们写出 $\Gamma$、$\Delta$、$\Pi$、$\Lambda$ 或 $\Theta$ 时，我们意指它们是我们所研究的逻辑中的公式多重集。按传统，证明论学者也不使用 $\tuple{\Gamma, \Delta}$ 的写法，而是用相继式符号分隔这两个部分；本书使用 $\Sequent$。

\begin{defn}[相继式]
一个\emph{相继式}是形如
\[
\Gamma \Sequent \Delta
\]
的表达式，其中 $\Gamma$ 和 $\Delta$ 是公式的有限（可能为空）多重集。$\Gamma$ 称为\emph{前件}，而 $\Delta$ 称为\emph{后件}。
\end{defn}

令 $!G$ 为 $\Gamma$ 中公式的合取（将空合取视为 $\ltrue$），$!D$ 为 $\Delta$ 中公式的析取（将空析取视为 $\lfalse$）。那么，相继式 $\Gamma \Sequent \Delta$（在某个结构 $\Struct{M}$ 中）为真，当且仅当 $!G \lif !D$ 为真。在经典逻辑中，这等价于说：要么 $\Gamma$ 中有一个公式为假，要么 $\Delta$ 中有一个公式为真。

如果 $\Gamma$ 是公式多重集，我们用 $\Gamma, !A$（或 $!A, \Gamma$）表示将 $!A$ 加入 $\Gamma$ 后得到的结果。如果 $\Delta$ 也是公式多重集，那么 $\Gamma, \Delta$ 就是这两个多重集的并——如果 $\Gamma$ 以\emph{重数} $n$ 包含 $!A$（即包含 $n$ 个 $!A$），而 $\Delta$ 以重数 $m$ 包含 $!A$，那么 $\Gamma, \Delta$ 就以重数 $n+m$ 包含 $!A$。如果相继式某一侧的多重集为空，我们通常就留出空白。例如，$\Sequent !A$ 就是一个前件为空、后件以重数 1 包含 $!A$ 的相继式。

相继式演算中的推导是由相继式构成的树，其中最顶端（或初始）的相继式是所考虑系统中的公理，并且若某个相继式位于一个或两个其他相继式的下方，则它必须是由系统中的推理规则正确推出的。我们将首先考虑经典逻辑的两个不同的相继式系统，$\Log{G1c}$ 和 $\Log{G3c}$。它们的公理和规则在 \cref{pt:seq:int:tab:G1c,pt:seq:int:tab:G3c} 中给出。\oliflabeldef{fol:seq::chap}{ （Gentzen 最初的系统 \Log{LK} 在 \olref[fol][seq][]{chap} 中描述。）}{}

\subfile{rules-G1c}
\subfile{rules-G3c}

这两个系统有细微差别，这些差别意味着它们具有不同的证明论性质。$\Log{G1c}$ 的公理为 $!A \Sequent !A$，且不允许包含其他公式。$\Log{G3c}$ 的公理形式为 $!C, \Gamma \Sequent \Delta, !C$，其中 $\Gamma$ 和 $\Delta$ 中可以有任意的“侧公式”，但两边出现的公式 $!C$ 必须是\emph{原子公式}。$\Log{G1c}$ 的 $\LeftR{\land}$ 规则和 $\RightR{\lor}$ 规则各有两个版本，而 $\Log{G3c}$ 各自只有一个版本。并且 $\Log{G1c}$ 拥有额外的“结构规则”，称为“收缩”（\LeftR{\Contraction}, \RightR{\Contraction}）和“弱化”（\LeftR{\Weakening}, \RightR{\Weakening}）。

系统 $\Log{G1c}$ 和 $\Log{G3c}$ 对经典逻辑恰好是可靠且完备的。换言之，在这些系统中可推导的每个相继式在每个结构中都为真，并且在每个结构中都为真的每个相继式都有一个推导。因此，一个相继式\emph{是否}有推导的问题，也可以通过逻辑的其他方法（例如模型论的方法）来回答。而且，由于两个系统恰好推导出那些在所有结构中都为真的相继式，特别地，它们证明相同的相继式。

然而，证明论不仅关心某物\emph{是否}可证，还关心它是\emph{如何}可证的。证明论学者会考虑诸如“能否总是避免使用某个特定规则？”、“能否将某条规则加入系统而不导致新的相继式变得可证？”、或者“能否将一个系统中的推导转变为另一个系统中的推导？”等问题。如果这类问题的答案是肯定的，那么对这一事实的证明通常通过对推导的复杂度进行归纳，或者等价地，对推导的结构进行归纳来完成。这不仅证明了该事实（例如，如果 \Log{G1c} 证明了某个相继式，那么 \Log{G3c} 也能证明同一个相继式），还提供了一个程序（例如，将 \Log{G1c} 中的推导转化为 \Log{G3c} 中的推导的程序）。

证明论的一个核心结果是\emph{切割消去定理}。它表明，我们可以将切割规则
\[
\Axiom$\Gamma \fCenter \Delta, !A$
\Axiom$!A, \Pi \fCenter \Lambda$
\RightLabel{\Cut}
\BinaryInf$\Gamma, \Pi \fCenter \Delta, \Lambda$
\DisplayProof
\]
添加到相继式系统中，而不会改变哪些相继式是可推导的。此外，该定理的证明提供了一种方法，能将任何同时使用了 \Log{G1c}（或 \Log{G3c}）规则与 \Cut{} 规则的推导，转换为只使用 \Log{G1c}（或 \Log{G3c}）规则的推导。换言之，该程序将 \Cut{} 规则从使用了它的推导中“消去”了；定理名便由此而来。

\end{document}